\lecture{2}
\title{Divide and Conquer}

\begin{document}

\textbf{Divide and Conquer.} Given any problem of size $N$, we\dots
\begin{itemize}
    \item \textbf{Divide.} Split the problem into $A$ subproblems, each of size $\frac{N}{B}$.
    \item \textbf{Conquer.} Solve each subproblem recursively.
    \item \textbf{Combine.} Use the subproblem solutions to get a solution of the original problem.
\end{itemize}
The runtime $T(N)$ is recursive: $T(N) = A \cdot T\left(\frac{N}{B}\right) + \text{(time to divide)} + \text{(time to combine)}$.

\begin{problem}{Rank Finding}
    Let $A$ be an unsorted array of $N$ numbers, and let $r \in [1, N]$ be a rank.
    
    Output the $r^{\text{th}}$ smallest number (i.e. the unique number with rank $r$).
\end{problem}

\textbf{Example.} If $r = 1$ or $r = N$, this asks for the minimum / maximum.  If $r = \frac{N + 1}{2}$, this asks for the median.

\begin{solution}{Rank Finding \#1}
    For each $x \in A$, compute $L(x) := \{ y \in A \mid y < x\}$.  Find the $x$ with $|L(x)| = r - 1$.
\end{solution}

Unfortunately, finding $L(x)$ takes $O(N)$ time, and so this is an $O(N^2)$ algorithm.  But this is also extremely wasteful:

\begin{boxed-text}[32em]
    \emph{Observation.} If $|L(x)| \geq r$, then our target must be in $L(x)$.
\end{boxed-text}

\begin{solution}{Rank Finding \#2}
    Pick a pivot $x \in A$, and consider $L(x) := \{ y \in A \mid y < x\}$ and $G(x) := \{ y \in A \mid y > x \}$.  Then recurse either on $L(x)$ or $G(x)$, depending on whether $|L(x)| \geq r$.
\end{solution}

Unfortunately, picking the pivot is important.  If we unluckily always pick $x = \min(A)$, this is $O(N^2)$ worst-case.

\begin{boxed-text}[32em]
    \emph{Observation.} If an oracle could tell us the median in $O(N)$ time, \\ then we'd have $T(N) = T\left(\frac{N}{2}\right) + O(N)$, which implies $T(N) = O(N)$.
\end{boxed-text}

In fact, we don't even need to pick \emph{exactly} the median.

\textbf{Definition.} For any $\frac{1}{2} \leq c < 1$, we say $x \in A$ is $c$-balanced if $\max \{ |L(x)|, |G(x)|\} \leq cN$.

\begin{boxed-text}[32em]
    \emph{Observation.} If an oracle could tell us a $c$-balanced element in $O(N)$ time, \\ then we'd have $T(N) = T\left(cN\right) + O(N)$, which implies $T(N) = O(N)$.
\end{boxed-text}

How do we find a $c$-balanced element?  The trick is to use Rank Finding again!

\begin{solution}{Rank Finding \#3}
    Divide $A$ into $\frac{N}{5}$ buckets of size $5$.  Say $M$ is the set of all bucket medians, with $|M| = \frac{N}{5}$.  Then set $x = \textsc{RankFindFast}\left(M, \frac{|M|}{2}\right)$ as the pivot, and proceed from there.
\end{solution}

\begin{proof}
    For algorithm correctness, it suffices to show:
    \begin{quote}
        \textbf{Claim.} Our choice of pivot $x$ is $\frac{7}{10}$-balanced.

        \emph{Proof:} The constant $\frac{7}{10}$ is $1 - \frac{1}{2} \times \frac{3}{5}$.  More generally, for buckets of size $2k + 1$, we have $c = \frac{3k + 1}{4k + 2}$.  ~ $\square$
    \end{quote}
    
    Now let's discuss the time complexity.  The recursion for this problem looks like:
    \[ T(N) = O(N) + T\left(\frac{N}{5}\right) + O(N) + T\left(\frac{7N}{10}\right) = T\left(\frac{N}{5}\right) + T\left(\frac{7N}{10}\right) + c_1 \cdot N, \]
    corresponding to (i) constructing $M$, (ii) finding $x$, (iii) computing $L(x)$ and $G(x)$, and (iv) solving the subproblem.

    Now we show by induction that $T(N) \leq c_2N$ for some $c_2 > 0$.  The inductive step is:
    \[ T(N) = T\left(\frac{N}{5}\right) + T\left(\frac{7N}{10}\right) + c_1 \cdot N \leq c_2 \cdot \left(\frac{9N}{10}\right) + c_1 \cdot N. \]
    And so of course $c_2 = 10c_1$ works, for example. ~ $\blacksquare$
\end{proof}

\begin{remark}
    The key success of this algorithm is that $\frac{1}{2k + 1} + \frac{3k + 1}{4k + 2} < 1$ when $k \geq 2$, meaning our algorithm does meaningfully cut down our search space by a factor.

    The reason why $k = 2$ is better than $k = 50$ is in step (i) constructing $M$.  Even though $\frac{1}{2k + 1} + \frac{3k + 1}{4k + 2}$ tends to $\frac{3}{4}$ for very large $k$, the constant $c_1$ also grows very large for large $k$.
\end{remark}

\begin{problem}{Integer Multiplication}
    Given two $N$-bit numbers $a$ and $b$, compute their product $a \cdot b$.
\end{problem}

\begin{solution}{Karatsuba}
    Write $a = y + 2^{N/2} \cdot x$ and $b = z + 2^{N/2} \cdot w$, so $a = [x] [y]$ and $b = [w] [z]$ and:
    \begin{align*}
        ab = \ & 2^N \cdot xw + 2^{N / 2} \cdot (xz + yw) + yz \\
        = \ & (2^N - 2^{N / 2}) \cdot xw + 2^{N / 2} \cdot (x + y)(w + z) + (1 - 2^{N / 2}) \cdot yz.
    \end{align*}
    So we only need $xw$, $yz$, and $(x + y)(w + z)$, which yields $T(N) = 3T\left(\frac{N}{2}\right) + O(N)$, or $T(N) = O(N^{\log_2(3)})$.
\end{solution}

In 1971, the Fast Fourier Transform was used to find an $O(N \cdot \log(N) \cdot \log(\log(N)))$ solution.

In 2007, $O(N \cdot \log(N) \cdot 2^{\log^*(N)})$ was found, where $\log^*(N)$ is the \# of $\log$'s with $\log(\log(\dots(\log(N))\dots)) < 1$.

And then in 2019, an $O(N \cdot \log(N))$ solution was found.  And we still don't know if we can do better.

\end{document}